% !PRINTED 398
\weberpage{398}
%<<CONT
reelle Doppeltangenten $z_3, z_4$ geben, so dass der Complex 4) sich
so fortsetzt:
\[
4)\quad x_1z_1,\ x_2z_2,\ x_3z_3,\ x_4z_4;
\]
$z_3$ und $z_4$ kommen nicht in 2) und auch nicht in 1) vor, und
müssen also in 3) enthalten sein.

Sie können auch in 3) nicht gepaart vorkommen, weil
$x_1, z_3, z_4$ azygetisch sind. Hieraus schliessen wir, wenn $t_3, t_4$ zwei
weitere reelle Doppeltangenten sind, auf folgende Zusammensetzung
 des Complexes 3):
\[
3)\quad x_1y_2,\ x_2y_1,\ z_3t_3,\ z_4t_4.
\]

Betrachtet man nun an Stelle des Complexes 4) den Complex
\[
5)\quad x_1z_3,\ y_2t_3,
\]
so kann man genau ebenso auf ein viertes reelles Paar $u_1u_2$ im
Complex 2) schliessen, und damit ist 5. bewiesen.

Fassen wir das Ergebniss zusammen, so erkennen wir, dass
in Bezug auf die Realität der Doppeltangenten einer reellen
Curve vierter Ordnung nur vier Fälle möglich sind:

1) \gesperrt{Vier reelle syzygetische Doppeltangenten.}

2) \gesperrt{Acht reelle Doppeltangenten, und zwar vier Paare
eines Steiner'schen Complexes.}

3) \gesperrt{Sechzehn reelle Doppeltangenten, die in einem
syzygetischen Complextripel je vier reelle Paare bilden.}

4) \gesperrt{Achtundzwanzig reelle Doppeltangenten.}

\begin{center}
\S.\,105.\\
\gesperrt{Beweis der Existenz der vier Fälle.}
\end{center}

Wir haben im vorigen Paragraphen zunächst nur bewiesen,
dass es keine anderen als die Fälle 1), 2), 3), 4) geben kann.
Dass diese vier Fälle aber wirklich alle möglich sind, ist jetzt
auch leicht zu zeigen, auf Grund des Satzes \S.\,100, nach dem
man aus sieben beliebig gegebenen geraden Linien auf \gesperrt{rationalem
 Wege} eine Curve vierter Ordnung ableiten kann, für
die die gegebenen Linien ein \gesperrt{Aronhold}'sches System bilden,
aus dem sich alle Doppeltangenten \gesperrt{rational} ableiten lassen.

Unter den sieben gegebenen geraden Linien können auch
imaginäre vorkommen, und als Bedingung der Realität der Curve
%CONT>>
